\documentclass[11pt,a4paper]{article}
\usepackage[margin=27mm]{geometry}
\usepackage{amsmath,amssymb,amsthm,hyperref,longtable,booktabs,array}
\usepackage[T1]{fontenc}
\usepackage{lmodern}
\emergencystretch=1em
\providecommand{\tightlist}{\setlength{\itemsep}{0pt}\setlength{\parskip}{0pt}}
\newcommand{\pandocbounded}[1]{#1}
\title{Pure subfield geometries and a scalar obstruction}
\author{Human Compression project}\date{}
\begin{document}\maketitle
For standard simple subfield parabolics, the exact criterion for
residual weak primitivity is that their common subfield remain simple.
Descent proves that the resulting geometries are thick, residually
connected and flag-transitive. In the mixed parabolic of Construction
10, maximality is instead equivalent to a scalar codimension of one;
this disproves the source claim for every prescribed parameter.

\section{Exact source scope and retained
inputs}\label{exact-source-scope-and-retained-inputs}

Write the odd field exponent as

\[
 n=\prod_{j=1}^{m}p_j^{f_j},\qquad p_i\ne p_j\ (i\ne j).
\]

Thus \(m\) counts distinct prime factors, whereas
\(\Omega(n)=\sum_j f_j\) counts factors with multiplicity. Construction
10, printed p.18, specifies rank \(3\le\alpha\le\min(m+1,\Omega(n))\),
selects \(k=\alpha-1\) pairwise nonisomorphic maximal Suzuki subfield
groups, puts \(a_i=n/\log_2(s_i)\), and sets \(t=2^{n/\prod_i a_i}\).
Its extra parabolic has type \((E_q\widehat{\cdot}E_t):C_{t-1}\).
Construction 11 instead specifies \(2\le\alpha\le m-1\) and selects
\(\alpha\) pairwise nonisomorphic maximal Suzuki subfield groups.
Theorem 5 immediately below claims that Constructions 1--11 give all
RWPRI pre-geometries of rank at least four. The rank bound and the
distinct-prime convention were checked in the page image, not only its
text extraction.

The source's definition on pp.4--5 is the exact condition used here: for
every proper subset \(J\) of the type set, at least one remaining
intersection \(G_{J\cup\{i\}}\) must be maximal in \(G_J\). In
particular, the condition applies to rank-one residues. The hat-dot on
p.5 means a nonsplit extension; \(E_q\) has order \(q\). The
introduction to the constructions on p.15 instructs the reader to retain
their whole intended intersection lattice, inferred from the rank-four
discussion. Construction 10 comes from the R12 discussion on pp.14--15.
Its intersections with the selected subfield groups have type
\((E_{s_i}\widehat{\cdot}E_t):C_{t-1}\).

We retain the following classical group theory.

\begin{enumerate}
\def\labelenumi{\arabic{enumi}.}
\tightlist
\item
  The Suzuki realization inside \(\operatorname{Sp}_4\), with
  exceptional Steinberg endomorphism \(\rho\), \(\rho^2=\phi_2\), and
  compatible standard groups
  \(S_d=\operatorname{Sz}(2^d)=\mathbf H^{\rho^d}\) for positive odd
  \(d\). Here \(S_1=\operatorname{Sz}(2)\cong5:4\) is allowed as an
  intersection, although it is not simple.
\item
  Prime-relative-degree inclusions \(S_d<S_{ad}\), for odd prime \(a\)
  and \(d>1\), are maximal. Standard subfield intersections are given by
  the gcd of their exponents.
\item
  The Suzuki Sylow-2 group \(U_q\) has exponent four and its identity
  together with its involutions is \(Z_q=Z(U_q)\), elementary abelian of
  order \(q\). A nonidentity involution \(z\in Z_q\) has \(C_G(z)=U_q\).
  Consequently a normalizer of \(Z_q\) normalizes \(U_q\), while a
  normalizer of \(U_q\) normalizes its center:
  \(N_G(Z_q)=N_G(U_q)=B_q\), the ovoid point stabilizer \(U_q:H_q\),
  where \(H_q\cong C_{q-1}\).
\item
  Nonsplit torus normalizers have order \(4(q\pm\sqrt{2q}+1)\) and are
  maximal for \(q\ge8\). Odd-element centralizers and these normalizers
  have their usual Suzuki structure.
\item
  Schur--Zassenhaus complement conjugacy for the solvable group
  \(U_q:H_t\). This puts its subgroups of order \(t-1\) into the
  standard \(H_t\).
\item
  Lang--Steinberg surjectivity on the connected algebraic group
  \(\operatorname{Sp}_4\), as used in the proved finite-descent and
  coset-Helly statements of Field descent and compatible tori, Section 2.
\end{enumerate}

The matrix action used in Section 4 is verified directly below. The
retained facts do not include any source classification of all RWPRI
lattices.

\section{Sharp criterion for the canonical pure-subfield
family}\label{sharp-criterion-for-the-canonical-pure-subfield-family}

\textbf{Theorem A.} Let \(n\) be positive and odd. Select \(k\ge2\)
distinct prime divisors \(a_1,\ldots,a_k\) of \(n\), put

\[
 e=\frac{n}{\prod_{i=1}^k a_i},\qquad
 G=S_n,\qquad G_i=S_{n/a_i}.
\]

The standard coset incidence system \(\Gamma(G;G_1,\ldots,G_k)\) is a
thick, residually connected, flag-transitive geometry. It is RWPRI
\textbf{if and only if \(e>1\)}. When \(e>1\), every one-step inclusion
in its intersection lattice, not only one inclusion at each subset, is
maximal.

\textbf{Proof.} For \(J\subseteq\{1,\ldots,k\}\), set

\[
 n_J=\frac{n}{\prod_{j\in J}a_j}.
\]

Entrywise field intersection gives \(G_J=S_{n_J}\), including
\(G_\varnothing=G\). Apply descent Theorem C with \(F=\rho^e\) and
\(m=\prod_i a_i\). Any pairwise incident family of cosets has a common
representative. It therefore extends to a translate of the standard
chamber. This proves the geometry axiom and flag transitivity.

Suppose that a residue has at least two remaining types. For distinct
remaining \(i,j\), the two residue parabolics \(S_{n_J/a_i}\) and
\(S_{n_J/a_j}\) are distinct maximal subgroups of \(S_{n_J}\): their
exponents exceed one, even if \(e=1\), because at least one other
relative prime remains. They generate \(S_{n_J}\). The standard
coset-graph connectivity criterion therefore proves that residue
connected. Flag transitivity identifies every residue with such a
standard one, proving RC.

The panel of type \(i\) has \([S_{ea_i}:S_e]\) chambers. This is at
least three. Indeed, \(S_{ea_i}\) is nonabelian simple and \(S_e\) is a
proper nontrivial subgroup, so its index cannot be two. This includes
\(e=1\), where \(|S_e|=20\). Thus the geometry is thick.

If \(e>1\), each remaining exponent \(n_J/a_i\) is at least \(e>1\).
Every one-step inclusion is a prime-degree inclusion of simple subfield
groups, and hence maximal. This proves RWPRI.

If \(e=1\), the distinct primes \(a_i\), with \(k\ge2\), include at
least one \(a_i\ge5\). At the subset containing every type except \(i\),
the only remaining inclusion is

\[
 S_1<S_{a_i}.
\]

Let \(C_5\) be the normal subgroup of \(S_1=5:4\). Inside \(S_{a_i}\),
its centralizer is a nonsplit torus. The order-four complement of
\(S_1\) normalizes this centralizer, so \(S_1\) lies inside the full
torus normalizer. Its torus order is one of
\(2^{a_i}\pm2^{(a_i+1)/2}+1\), both of which exceed five for
\(a_i\ge5\). Hence this containment is strict, while the torus
normalizer is proper in \(S_{a_i}\). Consequently \(S_1\) is not
maximal, and RWPRI fails at that rank-one residue. \(\square\)

The theorem concerns the standard simultaneous embeddings. It does not
assert that arbitrary conjugates of the same maximal-subgroup types have
the same intersections or are flag transitive. If another family has a
common prescribed simple subgroup \(S_e\) with \(e>1\), however, the
retained split-dihedral normalizer/transporter argument in
Earlier subfield and scalar proofs, Section B, forces simultaneous standard
position after conjugating \(S_e\). Thus within that prescribed
common-subfield lattice there is only one ambient conjugacy class for a
fixed labelled selection of exponents.

\section{Construction 11 and its omitted
boundary}\label{construction-11-and-its-omitted-boundary}

Construction 11 has \(k\le m-1\), so an unselected prime remains in
\(e\). In particular \(e>1\), and Theorem A supplies canonical
realizations for \textbf{every literal admissible parameter}. It also
answers the source's separate flag-transitivity question for this
family. It does not establish flag transitivity of all arbitrary
relative positions sharing only the maximal-parabolic isomorphism types.

The exact corrected parameter condition for the canonical pure-subfield
family is

\[
 2\le k\le m,\qquad \frac{n}{\prod_{i=1}^k a_i}>1.
\tag{1}
\]

For squarefree \(n\), this recovers \(k\le m-1\). If any prime factor of
\(n\) is repeated, it also permits \(k=m\). The source's bound omits
this latter boundary.

For example take

\[
 n=3^2\cdot5\cdot7\cdot11=3465,\qquad
 (a_1,a_2,a_3,a_4)=(3,5,7,11).
\]

Then

\[
 (G_1,G_2,G_3,G_4)
 =(S_{1155},S_{693},S_{495},S_{315}),\qquad
 G_{1234}=S_3=\operatorname{Sz}(8).
\]

Theorem A proves this is a rank-four thick, RC, FT, RWPRI geometry,
without enumerating its enormous ambient group. Here \(m=4\), whereas
Construction 11 allows only ranks two and three. Constructions 1--10
each have at least one soluble maximal parabolic, while all four
parabolics here are nonabelian simple. Consequently this stated-action
subgroup family is absent from all eleven literal constructions.

This contradicts the \textbf{literal exhaustive list} in Theorem 5, as a
theorem about the subgroup systems specified in the paper. It is a
boundary omission, distinct from proving that no natural extension of
the constructions could produce the example: allowing appropriate
truncations of larger constructions or replacing Construction 11's bound
by (1) would include it. The statement does not classify all unlabelled
geometries or all possible actions of \(G\).

\section{A scalar-center normal
form}\label{a-scalar-center-normal-form}

Here \(n\) and \(e\) are odd, \(e\mid n\), \(e>1\), \(q=2^n\), and
\(t=2^e\). Use standard Suzuki unipotents \(U(a,b)\), with
\(Z_q=\{U(0,b):b\in\mathbb F_q\}\), and diagonal matrices

\[
 h(\lambda)=
 \operatorname{diag}(\lambda^{\theta+1},\lambda,
                     \lambda^{-1},\lambda^{-\theta-1}),
 \qquad \theta=2^{(n+1)/2}.
\]

Multiplying these diagonal and unipotent matrices gives

\[
 h(\lambda)^{-1}U(a,b)h(\lambda)
 =U(\lambda^\theta a,\lambda^{\theta+2}b).
\tag{2}
\]

Both exponents are coprime to \(q-1\). For \(\theta+2\), if an odd
divisor \(c\) divided both \(q-1\) and \(2^{(n-1)/2}+1\), squaring the
latter congruence would give \(2^{n-1}\equiv1\pmod c\), whereas
\(2^n\equiv1\pmod c\); thus \(2\equiv1\pmod c\), forcing \(c=1\). Since
\(\theta+2=2(2^{(n-1)/2}+1)\), the claim follows. In particular the
standard \(H_t\) acts on both the first unipotent parameter and the
center by all nonzero scalars in \(\mathbb F_t\). Any additive invariant
subgroup is therefore an \(\mathbb F_t\)-subspace.

\textbf{Lemma B: normal form.} Every subgroup \(Q\le S_n\) of type
\((E_q\widehat{\cdot}E_t):C_{t-1}\) is ambient-conjugate to

\[
 Q_0=Z_qU_tH_t.
\tag{3}
\]

\textbf{Proof.} Its normal Sylow-2 subgroup \(P\) has order \(qt\) and
contains the elementary abelian kernel of order \(q\). Place \(P\)
inside an ambient \(U_q\). All its involutions lie in \(Z_q\), so that
elementary abelian kernel must equal \(Z_q\). Moreover \(Z_q\) is
exactly the identity together with the involutions of \(P\); it is
characteristic in \(P\) and normal in \(Q\). Therefore
\(Q\le N_G(Z_q)=B_q\).

Project to \(B_q/U_q=H_q\). The odd complement of \(Q\) maps
isomorphically to the unique subgroup of order \(t-1\), namely \(H_t\).
Complement conjugacy in \(U_q:H_t\) allows an ambient
\(U_q\)-conjugation putting it equal to \(H_t\). Now \(P/Z_q\) is an
additive subgroup \(W\) of the first unipotent parameter, of order
\(t\), normalized by \(H_t\). By (2), \(W=v\mathbb F_t\) for some
nonzero \(v\). A diagonal conjugation, which centralizes \(H_t\), scales
this line to \(\mathbb F_t\). Because \(P\) contains the full center, it
is the full preimage of that line in \(U_q\); thus \(P=Z_qU_t\). This
proves (3). \(\square\)

\textbf{Lemma C: exact maximality test.} Suppose \(H<Q_0\) has odd part
\(t-1\) and a normal Sylow-2 subgroup which is not elementary abelian.
Then \(H\) surjects onto \(Q_0/Z_q\). Put \(W=H\cap Z_q\), regarded as
an \(\mathbb F_t\)-subspace of \(\mathbb F_q\). The overgroups of \(H\)
in \(Q_0\) correspond exactly to the \(\mathbb F_t\)-subspaces
containing \(W\), by

\[
 V\longmapsto Z_VH,\qquad
 L\longmapsto L\cap Z_q.
\tag{4}
\]

In particular, \(H\) is maximal in \(Q_0\) exactly when
\(\dim_{\mathbb F_t}(\mathbb F_q/W)=1\).

\textbf{Proof.} A non-elementary Sylow-2 subgroup inside \(U_q\)
contains an element of order four. Its image modulo \(Z_q\) is nonzero,
so the image of \(H\) in \(Q_0/Z_q\) has a nonzero normal additive
subgroup. The odd projection has full order \(t-1\) and acts irreducibly
on the one-dimensional \(\mathbb F_t\)-space \(U_t/Z_t\). It follows
that the image contains all this additive group, and therefore equals
\(Q_0/Z_q\).

Thus \(Q_0=Z_qH\). The action of \(H\) on \(Z_q\) is exactly the scalar
action of \(H_t\), so \(W\) and every center intersection \(L\cap Z_q\)
are \(\mathbb F_t\)-subspaces. Conversely such a subspace \(V\) is
normalized by \(H\), making \(Z_VH\) a subgroup. Writing every element
as \(zh\) proves \(L=(L\cap Z_q)H\). This gives (4), with the stated
maximality criterion. \(\square\)

This is an overgroup classification inside the specified mixed
parabolic. It reduces maximality to one scalar-module dimension; no
ambient Suzuki group enumeration or search over relative positions is
required.

\section{Uniform obstruction for Construction
10}\label{uniform-obstruction-for-construction-10}

\textbf{Theorem D.} For every parameter allowed by Construction 10, the
intended bare subgroup-intersection lattice has a canonical embedding.
No embedding having its prescribed mixed parabolic and pairwise mixed
intersections is RWPRI.

\textbf{Proof.} Put \(k=\alpha-1\ge2\), select distinct odd relative
primes \(a_i\), and set

\[
 P=\prod_i a_i,\quad n=eP,\quad
 t=2^e,\quad q=2^n,\quad s_i=2^{n/a_i}.
\]

The source's bound \(\alpha\le\Omega(n)\) gives
\(\Omega(e)=\Omega(n)-k\ge1\), hence \(e>1\). Choose canonical subfield
groups \(G_i=S_{n/a_i}\) and \(Q=Z_qU_tH_t\). For every selection \(J\)
of subfield types,

\[
 \bigcap_{j\in J}G_j=S_{n_J},\qquad
 Q\cap\bigcap_{j\in J}G_j=Z_{2^{n_J}}U_tH_t.
\tag{5}
\]

These are exactly the intended field intersections and mixed-extension
orders, including common Borel \(U_tH_t\). The extension
\(Z_{2^{n_J}}U_t\) by its elementary abelian center kernel is nonsplit:
every element outside that center has order four, so a nontrivial
elementary abelian complement could not exist. Thus (5) also has the
literal hat-dot type. This proves bare existence.

For any other proposed embedding, Lemma B puts its mixed parabolic into
\(Q_0\). Its pairwise intersections \(H_i\) have the prescribed type
\((E_{s_i}\widehat{\cdot}E_t):C_{t-1}\). Their Sylow-2 groups cannot be
elementary abelian: an extension of elementary abelian groups whose
total group is elementary abelian splits as a vector-space extension.
Lemma C therefore applies, and order division gives \[
 |H_i\cap Z_q|=s_i,\qquad
 \dim_{\mathbb F_t}(\mathbb F_q/(H_i\cap Z_q))
 =P-P/a_i.
\]

Since \(k\ge2\) and every \(a_i\ge3\), this dimension exceeds one. Every
\(H_i\) is therefore nonmaximal in \(Q_0\). Explicitly, extend
\(H_i\cap Z_q\) by one \(\mathbb F_t\)-dimension; (4) gives a proper
intermediate subgroup. Taking the source's RWPRI subset \(J\) consisting
only of the mixed type leaves no maximal remaining intersection. This
contradicts RWPRI. \(\square\)

The negative conclusion uses only the mixed parabolic and the specified
pairwise intersection types. It does not require simultaneous alignment
of a common simple subfield. It strengthens the earlier rank-three R12
argument in Earlier subfield and scalar proofs, Section C. It does not prove
that arbitrary choices of the listed maximal parabolic types, with
different intersections, cannot satisfy RWPRI. Nor does failure of RWPRI
imply failure of the geometry or FT axioms.

\end{document}
